\section{基、坐标与维数：用最少的方向完整描述空间}
\label{sec:03-Linear-Algebra/02-Vector-Spaces/02}

你手里有三张地图，画的都是同一座城市。第一张画了南北方向，第二张画了东西方向，第三张把南北和东西各取一半画了一条斜线。三张图放在一起能不能说清楚城市里的每一条路？能。但有没有多余的图？也有——斜线那张不过是前两张拼出来的，它没告诉你任何新方向。

这场景搬到向量里一模一样。取
\[
e_1=\begin{pmatrix}1\\0\end{pmatrix},
\quad
e_2=\begin{pmatrix}0\\1\end{pmatrix},
\quad
e_1+e_2=\begin{pmatrix}1\\1\end{pmatrix}
\]
三个向量张成 \(\R^2\)，但第三个向量没有增加任何新方向。如果我们允许这种冗余，同一个向量就会有好几种系数写出来；可如果方向太少，又覆盖不了整个空间。

基刚好卡在两者之间：够用，不多余。把一叠地图精简到最薄的一本，每页都提供无法由其他页拼出来的方向——这就是基在做的事。

\subsection{线性无关：凭什么说一个向量是``多余''的}

先要有一个办法判断某个向量是否可以从其他向量拼出来。更准确地说，要判断一组向量里有没有谁在搭便车。

给定向量 \(v_1,\ldots,v_k\)。如果存在不全为零的标量 \(c_1,\ldots,c_k\) 使得
\[
c_1v_1+\cdots+c_kv_k=0,
\]
那就麻烦了。比如 \(c_3\neq0\)，把含 \(c_3\) 的项移到等号另一边再除以 \(c_3\)，就得到
\[
v_3=-\frac{c_1}{c_3}v_1-\cdots-\frac{c_k}{c_3}v_k.
\]
\(v_3\) 是别人的线性组合。它带不来新方向。

反过来说，如果
\[
c_1v_1+\cdots+c_kv_k=0
\]
迫使所有 \(c_i\) 都是零，那么没有任何一个向量能表示为其他人的组合。每个向量都提供了独立的方向。这套向量就叫线性无关。

看个具体例子：
\[
v_1=\begin{pmatrix}1\\0\\1\end{pmatrix},
\quad
v_2=\begin{pmatrix}0\\1\\1\end{pmatrix},
\quad
v_3=\begin{pmatrix}1\\1\\2\end{pmatrix}.
\]
一眼扫过去，第三个向量的前两个分量分别是 \(1\) 和 \(1\)，加起来刚好是第三个分量 \(2\)——这不是巧合。\(v_3=v_1+v_2\)，所以 \(1\cdot v_1+1\cdot v_2+(-1)\cdot v_3=0\)，系数不全为零。线性相关。删除 \(v_3\)，剩下的 \(v_1,v_2\) 张成的空间和原来完全一样。

这个判断过程和消元有一一对应。把向量作为矩阵的列：
\[
A=\begin{pmatrix}v_1&\cdots&v_k\end{pmatrix},
\]
那么 \(Ac=0\) 有没有非零解，就是在问这些列是不是线性相关。消元过程中出现的自由变量，正是列之间存在依赖关系的信号——一个自由变量对应一个可以被其他列表示的列。

所以线性无关不是在排除某个向量本身有什么问题，而是在追问：这个向量在不在别人能生成的范围里。如果在，就可以删；如果每删一个都损失了原来触及不了的方向，那它就在这个集合里拥有不可替代的位置。

\subsection{基：同时满足两个相反的要求}

基的定义是两个条件同时成立。对子空间 \(V\)，向量组 \(B=(v_1,\ldots,v_k)\) 是一组基，当且仅当：

\begin{enumerate}
\tightlist
\item
  它们张成 \(V\)；
\item
  它们线性无关。
\end{enumerate}

两个条件朝不同方向使劲。张成保证每个 \(v\in V\) 至少有一种表示方式；线性无关保证最多有一种。合在一起：每个 \(v\in V\) 恰好只有一种表示方式。

证明唯一性只用一步：
\[
v=c_1v_1+\cdots+c_kv_k
=d_1v_1+\cdots+d_kv_k,
\]
两式相减：
\[
(c_1-d_1)v_1+\cdots+(c_k-d_k)v_k=0.
\]
线性无关迫使每个 \(c_i-d_i=0\)，于是 \(c_i=d_i\)。这件事的后果比证明本身重要得多：基把空间里的每个对象和唯一的一组数字绑在一起。空间从此可以用数字说话。

如果把基看作一套好问题，每个基向量就是一个提问方向，坐标就是对这些问题的回答。不同的问题会给出不同的数字，但被描述的对象本身不会因此改变——这点等会儿在坐标部分会看得更清楚。

\subsection{坐标：同一个对象，不同的数字}

空间中的向量是一个几何或代数对象，坐标是用某组基对这个对象的描述。两者不是一回事。只有在标准基下，我们才偷懒把它们写成一样的数字——这恰好是因为标准基给了一个最省事的翻译。

若
\[
v=c_1v_1+\cdots+c_kv_k,
\]
则列向量
\[
[v]_B=
\begin{pmatrix}c_1\\\vdots\\c_k\end{pmatrix}
\]
称为 \(v\) 在基 \(B\) 下的坐标。注意这里的记号：\([v]_B\) 是一组数字，\(v\) 是向量本身；丢了角标 \(B\)，数字就失去了含义。

多项式空间最能揭示这个区别。在 \(P_2\) 中取基
\[
B=(1,t,t^2).
\]
多项式
\[
p(t)=2-3t+5t^2
\]
的坐标是 \((2,-3,5)^\trans\)。现在换一组基：
\[
C=(1,t-1,(t-1)^2).
\]
展开 \((t-1)^2=t^2-2t+1\)，整理：
\[
p(t)=2-3t+5t^2
=2-3t+5[(t-1)^2+2t-1]
=2-3t+5(t-1)^2+10t-5
=-3+7t+5(t-1)^2.
\]
再用 \(t=(t-1)+1\) 换掉主项：
\[
p(t)=-3+7[(t-1)+1]+5(t-1)^2
=4+7(t-1)+5(t-1)^2.
\]
所以
\[
[p]_C=(4,7,5)^\trans.
\]
多项式还是那个多项式，它的图像、根、导函数全都没变。两套坐标 \((2,-3,5)^\trans\) 和 \((4,7,5)^\trans\) 数字完全不同，却描述的是同一个数学对象。

这就是坐标的实质：坐标是对对象的翻译，不是对象本身。选什么基，就是在选用什么语言来和这个对象打交道。有的语言让计算简单，有的语言让结构暴露，语言本身不创造也不改变对象。

\subsection{一个平面告诉我们坐标怎么用}

考虑
\[
S=\{(x,y,z)^\trans:x+y+z=0\}.
\]
上一节找到了两个向量张成它：
\[
v_1=\begin{pmatrix}1\\-1\\0\end{pmatrix},
\qquad
v_2=\begin{pmatrix}1\\0\\-1\end{pmatrix}.
\]
它们线性无关吗？设
\[
c_1v_1+c_2v_2=0,
\]
写出分量：
\[
\begin{pmatrix}c_1+c_2\\-c_1\\-c_2\end{pmatrix}
=\begin{pmatrix}0\\0\\0\end{pmatrix}.
\]
从第二个分量立刻得到 \(c_1=0\)，第三个分量给出 \(c_2=0\)。所以 \(B=(v_1,v_2)\) 是 \(S\) 的一组基。

现在给定
\[
w=\begin{pmatrix}3\\-1\\-2\end{pmatrix},
\]
它在 \(B\) 下的坐标是什么？设
\[
c_1v_1+c_2v_2=w,
\]
分量：
\[
\begin{pmatrix}c_1+c_2\\-c_1\\-c_2\end{pmatrix}
=\begin{pmatrix}3\\-1\\-2\end{pmatrix}.
\]
从第二、三分量直接读出 \(c_1=1,c_2=2\)。验证第一分量 \(1+2=3\)，无误。于是
\[
[w]_B=\begin{pmatrix}1\\2\end{pmatrix}.
\]

注意一个有意思的事实：\(w\) 本来住在三维空间 \(\R^3\) 里，但它的坐标只有两个数。平面只有两个独立方向，所以只需要两个坐标。维度——也就是基向量的个数——直接反映了空间真正的自由度。

\subsection{维数：为什么它不依赖任何一组具体的基}

如果一个空间有两组基，它们的向量个数会不会不同？不会。

支撑这个结论的核心事实是：若 \(V\) 能由 \(n\) 个向量张成，那么 \(V\) 中任何线性无关组至多含 \(n\) 个向量。直觉很直白——\(n\) 个方向生成的房间，装不下超过 \(n\) 个互相独立的方向。严格证明通常写成替换引理：用线性无关组的向量逐个替换生成集中的成员，每次替换保持张成不变，最后比较替换步数。

有了这个事实，任取两组基：把第二组基看作生成集，第一组看作线性无关组，得到 \(|B_1|\le|B_2|\)；反过来交换角色得到 \(|B_2|\le|B_1|\)。所以任意两组基长度相同。这个共同的长度就是维数，记作
\[
\dim V.
\]

维数是一个空间最底层的数字签名。不管你怎么选基，它都不变：
\[
\dim\R^n=n,
\]
因为标准基 \(e_1,\ldots,e_n\) 有 \(n\) 个向量；
\[
\dim P_k=k+1,
\]
因为 \((1,t,\ldots,t^k)\) 是一组基，一共 \(k+1\) 个独立方向；所有 \(m\times n\) 实矩阵组成的空间维数为 \(mn\)——每个元素位置都是一个独立方向，取只有一个位置是 \(1\)、其余全为 \(0\) 的矩阵就得到一组基。

维数回答了这样一个问题：要完整描述这个空间里的任意对象，最少需要多少个独立参数。这是空间本身的性质，不是任何描述工具的产物。

\subsection{两个相反的操作：从生成集提取基，从无关组扩充基}

面对一组向量，你可能遇到两种情况。

情况一：它们张成了 \(V\)，但里面藏着线性相关。可以从冗余的成员开始逐个删：判断一个向量能否用剩下的向量表示，能就删。一直删到剩下的向量仍然张成 \(V\)，却不再包含任何冗余。得到一组基。

情况二：它们线性无关，但还没盖满 \(V\)。往里面加新向量：找一个不在当前张成空间中的向量塞进去，线性无关性保持，张成空间变大。一直加到张成 \(V\) 为止。也得到一组基。

两种操作互为镜像：
\begin{itemize}
\tightlist
\item
  生成集可以压缩为基；
\item
  线性无关组可以扩充为基。
\end{itemize}

这两种操作不只是一个存在性证明，它们给了我们一个实在的信念：任何有限生成向量空间都拥有一组基。基是这个空间的坐标骨架，每个这样的空间都能被一组基完整地钉住。

\subsection{消元如何帮你找基}

前面提到把向量作为矩阵 \(A\) 的列并做消元，主元列对应独立方向。但这个说法需要细化，因为行操作会同时改变每一列的具体数值。

行操作保持行空间。把行看成向量，初等行变换只是对行做线性组合，所以不改变行的张成空间。因此 RREF 的非零行可以直接作为原矩阵行空间的一组基。

但行操作改变列空间中的具体向量——每一列的具体数字都被重新组合了。所以列空间的基要从原矩阵的列中选，选与主元位置对应的那些列。主元位置告诉我们哪些列是独立的，而原矩阵中这些列的具体内容告诉我们这些独立方向落在哪个具体位置上。

二者的共同点是：主元数 \(r\) 同时是行秩和列秩。为什么行秩和列秩总相等？因为每个主元位置在行方向和列方向同时标记了一个独立维度。消元就像从矩阵中挤出了它的骨架，骨头数量从行方向数和从列方向数是一致的——同一个 \(r\)。

这个对应关系只依赖于行操作保持行空间、主元位置指示线性无关列这两条事实。到了正交单元，我们会看到这个 \(r\) 在四个基本子空间中的维数分布里扮演核心角色。

\subsection{维数限制：不需要算完就能判断的几条铁律}

在 \(n\) 维空间中，计数直接给出快速判断：

\begin{itemize}
\tightlist
\item
  超过 \(n\) 个向量必线性相关。房间只有 \(n\) 个独立方向，第 \(n+1\) 个方向无论如何都会落在前 \(n\) 个张成的空间中。
\item
  少于 \(n\) 个向量不可能张成整个空间。方向不够，必有余地够不到。
\item
  恰好 \(n\) 个向量时，线性无关和张成整个空间等价：证明了一个，另一个自动成立。
\end{itemize}

这三条怎么用？三个线性无关的 \(\R^3\) 向量自动构成一组基——不需要再验证它们是否张成 \(\R^3\)。同样，三个能张成 \(\R^3\) 的向量自动线性无关——不需要再检查系数。唯一需要确认的是数量对了、方向没有被浪费。

这些判断的威力在于：计数本身不告诉我们基是什么，但可以告诉我们什么时候该停下来，什么时候还差火候。面对一个 \(5\) 维空间拿到了六个向量，不必计算依赖关系就知道其中至少有一个是多余的；只拿出四个向量，不必检查张成范围就知道必然有方向落空。

\subsection{维数告诉你计算代价，但不告诉你怎么算}

知道一个空间的维数是 \(n\)，意味着用基表示任意向量需要 \(n\) 个数字。这 \(n\) 个数字在标准基下就是向量的分量本身；在一般的基下则需要解一个 \(n\times n\) 的线性系统来获得坐标。维数给出了复杂度的下界，但不同的基会让获得坐标的代价截然不同。下一节会看到，正交基把坐标计算从解方程组简化为点积，计算代价下降了一个量级。

维数也限制了矩阵能表达的信息量。一个秩为 \(r\) 的 \(m\times n\) 矩阵，输入有 \(n\) 维、输出有 \(m\) 维，但真正传递信息的通道只有 \(r\) 条。剩下的维度要么被零空间吞掉，要么是输出中永远达不到的方向。这套图景将在下一节``四个基本子空间''中完整展开。

延伸阅读：MIT 18.06：Independence, Basis and Dimension。
